
\documentclass[notes=show]{beamer}
%%%%%%%%%%%%%%%%%%%%%%%%%%%%%%%%%%%%%%%%%%%%%%%%%%%%%%%%%%%%%%%%%%%%%%%%%%%%%%%%%%%%%%%%%%%%%%%%%%%%%%%%%%%%%%%%%%%%%%%%%%%%%%%%%%%%%%%%%%%%%%%%%%%%%%%%%%%%%%%%%%%%%%%%%%%%%%%%%%%%%%%%%%%%%%%%%%%%%%%%%%%%%%%%%%%%%%%%%%%%%%%%%%%%%%%%%%%%%%%%%%%%%%%%%%%%
\usepackage{mathpazo}
\usepackage{hyperref}
\usepackage{multimedia}

%TCIDATA{OutputFilter=LATEX.DLL}
%TCIDATA{Version=5.50.0.2953}
%TCIDATA{<META NAME="SaveForMode" CONTENT="1">}
%TCIDATA{BibliographyScheme=Manual}
%TCIDATA{Created=Monday, September 29, 2008 11:10:36}
%TCIDATA{LastRevised=Tuesday, June 21, 2022 18:25:45}
%TCIDATA{<META NAME="GraphicsSave" CONTENT="32">}
%TCIDATA{<META NAME="DocumentShell" CONTENT="Other Documents\SW\Slides - Beamer">}
%TCIDATA{CSTFile=beamer.cst}

\newenvironment{stepenumerate}{\begin{enumerate}[<+->]}{\end{enumerate}}
\newenvironment{stepitemize}{\begin{itemize}[<+->]}{\end{itemize} }
\newenvironment{stepenumeratewithalert}{\begin{enumerate}[<+-| alert@+>]}{\end{enumerate}}
\newenvironment{stepitemizewithalert}{\begin{itemize}[<+-| alert@+>]}{\end{itemize} }
\usetheme{Madrid}
\input{tcilatex}
\begin{document}

\title{Macroeconomics II\\
Stochastic Recursive Models}
\author{Prof. McCandless}
\maketitle

\section{Stochastic Recursive Models}

%TCIMACRO{\TeXButton{BeginFrame}{\begin{frame}}}%
%BeginExpansion
\begin{frame}%
%EndExpansion

\QTR{frametitle}{Probability}

\begin{itemize}
\item A probability space is a triplet:%
\[
\left( \Omega ,\tciFourier ,P\right) 
\]

\item Where

\begin{itemize}
\item $\Omega $ is a set of all states of nature that can occur

\item $\tciFourier $ is a collection of subsets of $\Omega $, each subset is
called an "event"

\item $P$ is a probability measure over events, i.e. $\tciFourier $
\end{itemize}

\item Note that I said a probability \textbf{measure}. \ Measures tell you
how big or far apart (in some mathematical sense) points or sets are.

\item These can be used in integration, for example, let p(x) be the
probability of event x, in this case with elements of x taken from a line
between 0 and 10, then \ 
\[
\int_{0}^{10}p(x)dx=1
\]
\end{itemize}

%TCIMACRO{\TeXButton{EndFrame}{\end{frame}}}%
%BeginExpansion
\end{frame}%
%EndExpansion

%TCIMACRO{\TeXButton{BeginFrame}{\begin{frame}}}%
%BeginExpansion
\begin{frame}%
%EndExpansion

\QTR{frametitle}{Probability for finite states of nature}

\begin{itemize}
\item Suppose that $\Omega $ is a finite set

\item We can sometimes impose the probability measure directly on elements
of $\Omega $

\begin{itemize}
\item if each state of nature can be considered an event
\end{itemize}

\item Define $p_{i}=$ probability that event $A_{i}$ will occur

\item If the probabilities are independent ($A_{i}$ and $A_{j}$ do not
intersect, have no states in common)

\begin{itemize}
\item $p\left( A_{i},A_{j}\right) =p_{i}+p_{j}$
\end{itemize}

\item We can have non-independent events (where events share possible states)

\begin{itemize}
\item $E_{1}=\left( A_{1},A_{2}\right) $

\item $E_{2}=\left( A_{2},A_{3}\right) $

\item Then $p(E_{1}\cup E_{2})=p_{1}+p_{2}+p_{3}<p(E_{1})+p(E_{2})$
\end{itemize}
\end{itemize}

%TCIMACRO{\TeXButton{EndFrame}{\end{frame}}}%
%BeginExpansion
\end{frame}%
%EndExpansion

%TCIMACRO{\TeXButton{BeginFrame}{\begin{frame}}}%
%BeginExpansion
\begin{frame}%
%EndExpansion

\QTR{frametitle}{Probability for finite states of nature}

\begin{itemize}
\item Let $\Omega =\left\{ A_{1},A_{2}\right\} $

\item Then the largest possible $\tciFourier $ is

\begin{itemize}
\item $[]$, $[A_{1}]$, $[A_{2}]$, $[A_{1},A_{2}]$
\end{itemize}

\item A possible probability measure is

\begin{itemize}
\item $p\left( []\right) =0$, $p\left( [A_{1}]\right) =\overline{p}$, $%
p\left( [A_{2}]\right) =1-\overline{p}$, and $p\left( [A_{1},A_{2}]\right)
=1 $
\end{itemize}

\item Another possible is

\begin{itemize}
\item $p\left( []\right) =\widehat{p}_{1}$, $p\left( [A_{1}]\right) =%
\overline{p}$, $p\left( [A_{2}]\right) =1-\widehat{p}_{1}-\overline{p}$, and 
$p\left( [A_{1},A_{2}]\right) =1-\widehat{p}_{1}$
\end{itemize}
\end{itemize}

%TCIMACRO{\TeXButton{EndFrame}{\end{frame}}}%
%BeginExpansion
\end{frame}%
%EndExpansion

%TCIMACRO{\TeXButton{BeginFrame}{\begin{frame}}}%
%BeginExpansion
\begin{frame}%
%EndExpansion

\QTR{frametitle}{Probability for continuous states of nature}

\begin{itemize}
\item Let $\Omega =\left[ 0,300\right] $, all the points on the line from $0$
to $300$ (inclusive)

\item Let $\tciFourier $ be all measureable sets of $\Omega $

\item $P$ assigns probabilities to these measurable sets

\item $p(\left[ 45,46.1\right] )$ is the probability of the value falling
between $45$ and $46.1$

\item $p\left( \pi ,4\right) $ is the probability of the value falling
between $\pi $\ and $4$

\item In general, $p\left( x\right) =0$, where $x$ is a specific number

\begin{itemize}
\item Not always the case

\begin{itemize}
\item $p\left( x\right) $, where $x$ is the daily rainfall in Buenos Aries

\item $p\left( 0\right) $ is the probability it won't rain

\item $p\left( 0\right) >0$
\end{itemize}
\end{itemize}
\end{itemize}

%TCIMACRO{\TeXButton{EndFrame}{\end{frame}}}%
%BeginExpansion
\end{frame}%
%EndExpansion

%TCIMACRO{\TeXButton{BeginFrame}{\begin{frame}}}%
%BeginExpansion
\begin{frame}%
%EndExpansion

\QTR{frametitle}{A simple stochastic model}

\begin{itemize}
\item Robinson Crusoe model with stochastic technology

\item Production function%
\[
y_{t}=A^{t}f(k_{t}), 
\]%
with the technology $A^{t}$ determined by 
\[
A^{t}=\left\{ 
\begin{array}{c}
A_{1}\text{ with probability }p_{1} \\ 
A_{2}\text{ with probability }p_{2}%
\end{array}%
\right. 
\]%
and we assume that $A_{1}>A_{2}$ and that $p_{2}=1-p_{1}$

\item RC's utility function is%
\[
E_{0}\sum_{t=0}^{\infty }\beta ^{t}u(c_{t}) 
\]%
subject to the budget constraint%
\[
k_{t+1}=A^{t}f(k_{t})+(1-\delta )k_{t}-c_{t}. 
\]

\item This is an \textbf{expected} utility function: that is what $E_{0}$
means
\end{itemize}

%TCIMACRO{\TeXButton{EndFrame}{\end{frame}}}%
%BeginExpansion
\end{frame}%
%EndExpansion

%TCIMACRO{\TeXButton{BeginFrame}{\begin{frame}}}%
%BeginExpansion
\begin{frame}%
%EndExpansion

\QTR{frametitle}{Stochastic Value function}

\begin{itemize}
\item The value of discounted expected utility at time $0$ when the realized
technology shock is $A_{1}$\ is%
\[
V(k_{0},A_{1})=\max_{\left\{ c_{t}\right\} _{t=0}^{\infty
}}E_{0}\sum_{t=0}^{\infty }\beta ^{t}u(c_{t}),
\]%
subject to the budget constraint for $t=0$,%
\[
k_{1}=A_{1}f(k_{0})+(1-\delta )k_{0}-c_{0},
\]%
and those for $t\geq 1,$%
\[
k_{t+1}=A^{t}f(k_{t})+(1-\delta )k_{t}-c_{t},
\]%
and the independent realizations of $A^{t}=\left[ A_{1},A_{2}\right] $ with
probabilities $\left[ p_{1},p_{2}\right] $

\item There is a similar expression for $V(k_{0},A_{2})$
\end{itemize}

%TCIMACRO{\TeXButton{EndFrame}{\end{frame}}}%
%BeginExpansion
\end{frame}%
%EndExpansion

%TCIMACRO{\TeXButton{BeginFrame}{\begin{frame}}}%
%BeginExpansion
\begin{frame}%
%EndExpansion

\QTR{frametitle}{Stochastic Value function: recursive format}

\begin{itemize}
\item The value function is%
\[
V(k_{0},A^{0})=\max_{c_{0}}u(c_{0})+\beta E_{0}V\left( k_{1},A^{1}\right) 
\]%
subject to the budget constraint%
\[
k_{1}=A^{0}f(k_{0})+(1-\delta )k_{0}-c_{0}. 
\]

\item Here we are taking $c_{0}$ as the control

\item The states are $k_{0}$ and $A^{0}$

\item Notice the expectations operator

\begin{itemize}
\item In the second part of the value function

\item $E_{0}$ operator says we don't know the value of time 1 variables, but
we know their probabilities

\item we do not know the realization of $A^{1}$ but we know their possible
values
\end{itemize}
\end{itemize}

%TCIMACRO{\TeXButton{EndFrame}{\end{frame}}}%
%BeginExpansion
\end{frame}%
%EndExpansion

%TCIMACRO{\TeXButton{BeginFrame}{\begin{frame}}}%
%BeginExpansion
\begin{frame}%
%EndExpansion

\QTR{frametitle}{Stochastic Value function: with k(t+1)\ as control}

\begin{itemize}
\item The value function is%
\[
V(k_{t},A^{t})=\max_{k_{t+1}}u(A^{t}f(k_{t})+(1-\delta )k_{t}-k_{t+1})+\beta
E_{t}V\left( k_{t+1},A^{t+1}\right) 
\]%
and the budget constraint is%
\[
k_{t+1}=G(x_{t},y_{t})=k_{t+1} 
\]

\item The control at time $t$ is the state at time $t+1$

\item We solve for a \textbf{plan}, a function such that 
\[
k_{t+1}=H(k_{t},A^{t}) 
\]

\item A plan solves (without maximization)%
\begin{eqnarray*}
V(k_{t},A^{t}) &=&u(A^{t}f(k_{t})+(1-\delta )k_{t}-H(k_{t},A^{t})) \\
&&+\beta E_{t}V\left( H(k_{t},A^{t}),A^{t+1}\right)
\end{eqnarray*}
\end{itemize}

%TCIMACRO{\TeXButton{EndFrame}{\end{frame}}}%
%BeginExpansion
\end{frame}%
%EndExpansion

%TCIMACRO{\TeXButton{BeginFrame}{\begin{frame}}}%
%BeginExpansion
\begin{frame}%
%EndExpansion

\QTR{frametitle}{General version of the problem}

\begin{itemize}
\item Write the value function as%
\[
V(x_{t},z_{t})=\max_{\left\{ y_{s}\right\} _{s=t}^{\infty
}}E_{t}\sum_{s=t}^{\infty }\beta ^{s-t}F(x_{s},y_{s},z_{s}), 
\]%
subject to 
\[
x_{s+1}=G(x_{s},y_{s},z_{s})\text{for }s\geq t 
\]

\item $x_{t}$ is the set of "regular", predetermined, state variables

\item $z_{t}$ is the set of state variables determined by nature

\begin{itemize}
\item these are the stochastic state variables.
\end{itemize}

\item $y_{t}$ are the control variables \ 

\item Both $F(x_{s},y_{s},z_{s})$ and $G(x_{s},y_{s},z_{s})$ can contain the
stochastic state variables.
\end{itemize}

%TCIMACRO{\TeXButton{EndFrame}{\end{frame}}}%
%BeginExpansion
\end{frame}%
%EndExpansion

%TCIMACRO{\TeXButton{BeginFrame}{\begin{frame}}}%
%BeginExpansion
\begin{frame}%
%EndExpansion

\QTR{frametitle}{General version of the problem}

\begin{itemize}
\item The recrusive version of this problem is%
\[
V(x_{t},z_{t})=\max_{y_{t}}\left[ F(x_{t},y_{t},z_{t})+\beta
E_{t}V(x_{t+1},z_{t+1})\right] 
\]%
subject to 
\[
x_{t+1}=G(x_{t},y_{t},z_{t}) 
\]

\item A solution is a plan,%
\[
y_{t}=H(x_{t},z_{t}) 
\]%
where%
\begin{eqnarray*}
V(x_{t},z_{t}) &=&F(x_{t},H(x_{t},z_{t}),z_{t}) \\
&&+\beta E_{t}V(G(x_{t},H(x_{t},z_{t}),z_{t}),z_{t+1})
\end{eqnarray*}
\end{itemize}

%TCIMACRO{\TeXButton{EndFrame}{\end{frame}}}%
%BeginExpansion
\end{frame}%
%EndExpansion

%TCIMACRO{\TeXButton{BeginFrame}{\begin{frame}}}%
%BeginExpansion
\begin{frame}%
%EndExpansion

\QTR{frametitle}{General version of the problem: first order conditions}

\begin{itemize}
\item The first order conditions are%
\[
0=F_{y}(x_{t},y_{t},z_{t})+\beta E_{t}\left[
V_{x}(G(x_{t},y_{t},z_{t}),z_{t+1})G_{y}(x_{t},y_{t},z_{t})\right] 
\]

\item The Benveniste-Scheinkman condition is%
\[
V_{x}(x_{t},z_{t})=F_{x}(x_{t},y_{t},z_{t})+\beta E_{t}\left[
V_{x}(G(x_{t},y_{t},z_{t}),z_{t+1})G_{x}(x_{t},y_{t},z_{t})\right] 
\]

\item If we can choose the controls so that $G_{x}(x_{t},y_{t},z_{t})=0$,
this becomes%
\[
V_{x}(x_{t},z_{t})=F_{x}(x_{t},y_{t},z_{t}) 
\]

\item One can write the \textbf{stochastic Euler equation} as%
\[
0=F_{y}(x_{t},y_{t},z_{t})+\beta E_{t}\left[
F_{x}(G(x_{t},y_{t},z_{t}),y_{t+1},z_{t+1})G_{y}(x_{t},y_{t},z_{t})\right] 
\]
\end{itemize}

%TCIMACRO{\TeXButton{EndFrame}{\end{frame}}}%
%BeginExpansion
\end{frame}%
%EndExpansion

%TCIMACRO{\TeXButton{BeginFrame}{\begin{frame}}}%
%BeginExpansion
\begin{frame}%
%EndExpansion

\QTR{frametitle}{Solving for the value function}

\begin{itemize}
\item One can find an approximation of the value function from%
\[
V_{j+1}(x_{t},z_{t})=\max_{y_{t}}\left[ F(x_{t},y_{t},z_{t})+\beta
E_{t}V_{j}(G(x_{t},y_{t},z_{t}),z_{t+1})\right] 
\]

\item Beginning with some function $V_{0}(\cdot )$ (frequently a constant)

\item Need to solve over sufficiently dense set of $X\times Z$

\begin{itemize}
\item where $X$ is the domain of the state variables, $x_{t}$

\item $Z$ is the domain of the states of nature, $z_{t}$

\item sufficient depends on the level of precision we need
\end{itemize}
\end{itemize}

%TCIMACRO{\TeXButton{EndFrame}{\end{frame}}}%
%BeginExpansion
\end{frame}%
%EndExpansion

%TCIMACRO{\TeXButton{BeginFrame}{\begin{frame}}}%
%BeginExpansion
\begin{frame}%
%EndExpansion

\QTR{frametitle}{Problem of dimensionality}

\begin{itemize}
\item Problem of dimensionality is worse than in the deterministic case

\item In the deterministic case is based on

\begin{itemize}
\item the number of predetermined state variables

\item the size of the sufficiently dense subset of each state variable we use
\end{itemize}

\item In the stochastic state

\begin{itemize}
\item these two problems continue

\item add

\begin{itemize}
\item the dimension of the shocks (if finite)

\item the sufficiently dense subset of the shocks (if continuous)
\end{itemize}
\end{itemize}
\end{itemize}

%TCIMACRO{\TeXButton{EndFrame}{\end{frame}}}%
%BeginExpansion
\end{frame}%
%EndExpansion

%TCIMACRO{\TeXButton{BeginFrame}{\begin{frame}}}%
%BeginExpansion
\begin{frame}%
%EndExpansion

\QTR{frametitle}{Finding the value function for our simple economy}

\begin{itemize}
\item In the growth economy, the technology level can be $\left[ A_{1},A_{2}%
\right] $

\item This gives two values functions of the form%
\begin{eqnarray*}
V(k_{t},A_{1}) &=&\max_{k_{t+1}}u(A_{1}f(k_{t})+(1-\delta )k_{t}-k_{t+1}) \\
&&+\beta \left[ p_{1}V\left( k_{t+1},A_{1}\right) +p_{2}V\left(
k_{t+1},A_{2}\right) \right]
\end{eqnarray*}%
and%
\begin{eqnarray*}
V(k_{t},A_{2}) &=&\max_{k_{t+1}}u(A_{2}f(k_{t})+(1-\delta )k_{t}-k_{t+1}) \\
&&+\beta \left[ p_{1}V\left( k_{t+1},A_{1}\right) +p_{2}V\left(
k_{t+1},A_{2}\right) \right]
\end{eqnarray*}

\item Notice how the probabilities enter

\item because the shocks have oen of two values, we need to find two
functions, $V(\cdot ,A_{1})$ and $V(\cdot ,A_{2})$
\end{itemize}

%TCIMACRO{\TeXButton{EndFrame}{\end{frame}}}%
%BeginExpansion
\end{frame}%
%EndExpansion

%TCIMACRO{\TeXButton{BeginFrame}{\begin{frame}}}%
%BeginExpansion
\begin{frame}%
%EndExpansion

\QTR{frametitle}{The recursive approximation}

\begin{itemize}
\item Same recursive approximation as before (as in the deterministic
version)

\item Difference is that we need to find two equations at each iteration

\item Given $V_{0}(\cdot ,A_{1})$ and $V_{0}(\cdot ,A_{2})$, we find%
\begin{eqnarray*}
V_{1}(k_{t},A_{1}) &=&\max_{k_{t+1}}u(A_{1}f(k_{t})+(1-\delta )k_{t}-k_{t+1})
\\
&&+\beta \left[ p_{1}V_{0}\left( k_{t+1},A_{1}\right) +p_{2}V_{0}\left(
k_{t+1},A_{2}\right) \right] ,
\end{eqnarray*}%
and%
\begin{eqnarray*}
V_{1}(k_{t},A_{2}) &=&\max_{k_{t+1}}u(A_{2}f(k_{t})+(1-\delta )k_{t}-k_{t+1})
\\
&&+\beta \left[ p_{1}V_{0}\left( k_{t+1},A_{1}\right) +p_{2}V_{0}\left(
k_{t+1},A_{2}\right) \right] ,
\end{eqnarray*}

\item Repeat, finding $V_{N}(\cdot ,A_{1})$ and $V_{N}(\cdot ,A_{2})$ until
sufficiently close
\end{itemize}

%TCIMACRO{\TeXButton{EndFrame}{\end{frame}}}%
%BeginExpansion
\end{frame}%
%EndExpansion

%TCIMACRO{\TeXButton{BeginFrame}{\begin{frame}}}%
%BeginExpansion
\begin{frame}%
%EndExpansion

\QTR{frametitle}{Example}

\begin{itemize}
\item Used $\delta =.1$, $\beta =.98$, $A_{1}=1.75$, $p_{1}=.8$, $A_{2}=.75$,

\item and $p_{2}=.2$, $V_{0}(\cdot ,A_{1})=20$ and $V_{0}(\cdot ,A_{2})=20$

\item Graph of iterations\FRAME{ftbpFU}{3.039in}{2.2857in}{0pt}{\Qcb{%
Iterations on the value function}}{}{value fun sto.jpg}{\special{language
"Scientific Word";type "GRAPHIC";maintain-aspect-ratio TRUE;display
"USEDEF";valid_file "F";width 3.039in;height 2.2857in;depth
0pt;original-width 5.8332in;original-height 4.3751in;cropleft "0";croptop
"1";cropright "1";cropbottom "0";filename 'C:/macro 2 2022/class 3/value fun
sto.jpg';file-properties "XNPEU";}}
\end{itemize}

%TCIMACRO{\TeXButton{EndFrame}{\end{frame}}}%
%BeginExpansion
\end{frame}%
%EndExpansion

%TCIMACRO{\TeXButton{BeginFrame}{\begin{frame}}}%
%BeginExpansion
\begin{frame}%
%EndExpansion

\QTR{frametitle}{The two policy functions (the plans)}

\FRAME{ftbpFU}{3.039in}{2.2857in}{0pt}{\Qcb{The plans}}{}{policy function
sto.jpg}{\special{language "Scientific Word";type
"GRAPHIC";maintain-aspect-ratio TRUE;display "USEDEF";valid_file "F";width
3.039in;height 2.2857in;depth 0pt;original-width 5.8332in;original-height
4.3751in;cropleft "0";croptop "1";cropright "1";cropbottom "0";filename
'C:/macro 2 2022/class 3/policy function sto.jpg';file-properties "XNPEU";}}

%TCIMACRO{\TeXButton{EndFrame}{\end{frame}}}%
%BeginExpansion
\end{frame}%
%EndExpansion

%TCIMACRO{\TeXButton{BeginFrame}{\begin{frame}}}%
%BeginExpansion
\begin{frame}%
%EndExpansion

\QTR{frametitle}{A simulation of the economy}

\FRAME{ftbpFU}{3.039in}{2.2857in}{0pt}{\Qcb{A simulated time path}}{}{%
storun1 exp.jpg}{\special{language "Scientific Word";type
"GRAPHIC";maintain-aspect-ratio TRUE;display "USEDEF";valid_file "F";width
3.039in;height 2.2857in;depth 0pt;original-width 5.8332in;original-height
4.3751in;cropleft "0";croptop "1";cropright "1";cropbottom "0";filename
'C:/macro 2 2022/class 3/storun1 exp.jpg';file-properties "XNPEU";}}

%TCIMACRO{\TeXButton{EndFrame}{\end{frame}}}%
%BeginExpansion
\end{frame}%
%EndExpansion

%TCIMACRO{\TeXButton{BeginFrame}{\begin{frame}}}%
%BeginExpansion
\begin{frame}%
%EndExpansion

\QTR{frametitle}{Markov chains}

\begin{itemize}
\item The above simulation shows relatively little persistence

\item Markov chains are a way of adding persistence to the shocks

\begin{itemize}
\item Note that the presistence is in the stochastic variable

\item The \textbf{economic} model is not generating this persistence
\end{itemize}

\item Structure of a Markov chain: conditional probabilities

\begin{itemize}
\item The probabilities at time $t$ of the time $t+1$ states of nature

\item depend on the state of nature at time $t$
\end{itemize}

\item Consider our example with two states of nature $\left[ A_{1},A_{2}%
\right] $

\item Let the probabilities be%
\[
P=\left[ 
\begin{array}{cc}
p_{11} & p_{12} \\ 
p_{21} & p_{22}%
\end{array}%
\right] 
\]%
where $p_{ij}$ is the probability of going to state $j$ given you are in
state $i$
\end{itemize}

%TCIMACRO{\TeXButton{EndFrame}{\end{frame}}}%
%BeginExpansion
\end{frame}%
%EndExpansion

%TCIMACRO{\TeXButton{BeginFrame}{\begin{frame}}}%
%BeginExpansion
\begin{frame}%
%EndExpansion

Probabilities in Markov chains

\begin{itemize}
\item These are \textbf{conditional probabilities}

\item If one is in state of nature $1$ at time $t$

\begin{itemize}
\item The probabilities for time $t+1$ are 
\[
\left[ 
\begin{array}{cc}
p_{11} & p_{12}%
\end{array}%
\right] 
\]
\end{itemize}

\item If one is in state of nature $2$ at time $t$

\begin{itemize}
\item The probabilities for time $t+1$ are 
\[
\left[ 
\begin{array}{cc}
p_{21} & p_{22}%
\end{array}%
\right] 
\]
\end{itemize}

\item Example with a lot of persistence%
\[
P=\left[ 
\begin{array}{cc}
.97 & .03 \\ 
.1 & .9%
\end{array}%
\right] 
\]
\end{itemize}

%TCIMACRO{\TeXButton{EndFrame}{\end{frame}}}%
%BeginExpansion
\end{frame}%
%EndExpansion

%TCIMACRO{\TeXButton{BeginFrame}{\begin{frame}}}%
%BeginExpansion
\begin{frame}%
%EndExpansion

\QTR{frametitle}{Unconditional probabilites}

\begin{itemize}
\item What is the probability that one will be in state of nature $j$ at
some far distant date

\item Does this depend on the current state of nature

\item Given the state at time $0$, the distributiion for period $1$ is $%
p_{0}=\left[ 
\begin{array}{cc}
p_{01} & p_{02}%
\end{array}%
\right] $

\item Then the distribution for period $2$ is 
\[
p_{0}P=\left[ 
\begin{array}{cc}
p_{01} & p_{02}%
\end{array}%
\right] \left[ 
\begin{array}{cc}
p_{11} & p_{12} \\ 
p_{21} & p_{22}%
\end{array}%
\right] 
\]

\item The distribution for period $3$ is 
\[
p_{0}PP=\left[ 
\begin{array}{cc}
p_{01} & p_{02}%
\end{array}%
\right] \left[ 
\begin{array}{cc}
p_{11} & p_{12} \\ 
p_{21} & p_{22}%
\end{array}%
\right] \left[ 
\begin{array}{cc}
p_{11} & p_{12} \\ 
p_{21} & p_{22}%
\end{array}%
\right] 
\]

\item The distribution for period $N$ is%
\[
p_{0}P^{N-1}
\]
\end{itemize}

%TCIMACRO{\TeXButton{EndFrame}{\end{frame}}}%
%BeginExpansion
\end{frame}%
%EndExpansion

%TCIMACRO{\TeXButton{BeginFrame}{\begin{frame}}}%
%BeginExpansion
\begin{frame}%
%EndExpansion

\QTR{frametitle}{Converge to an unconditional probability}

\begin{itemize}
\item What happens as $N$ gets large

\item Use our example probability matrix (leave $p_{0}$ out for the moment)

\item Start with%
\[
P=\left[ 
\begin{array}{cc}
.97 & .03 \\ 
.1 & .9%
\end{array}%
\right] 
\]

\item $PP=P^{2}$ is

\item 
\[
P^{2}=\left[ 
\begin{array}{cc}
.97 & .03 \\ 
.1 & .9%
\end{array}%
\right] \left[ 
\begin{array}{cc}
.97 & .03 \\ 
.1 & .9%
\end{array}%
\right] =\left[ 
\begin{array}{cc}
0.9439 & 0.0561 \\ 
0.1870 & 0.8130%
\end{array}%
\right] 
\]
\end{itemize}

%TCIMACRO{\TeXButton{EndFrame}{\end{frame}}}%
%BeginExpansion
\end{frame}%
%EndExpansion

%TCIMACRO{\TeXButton{BeginFrame}{\begin{frame}}}%
%BeginExpansion
\begin{frame}%
%EndExpansion

\QTR{frametitle}{Using a doubling algorythm (which saves time)}%
\[
P^{4}=P^{2}P^{2}=\left[ 
\begin{array}{cc}
0.944 & 0.056 \\ 
0.187 & 0.813%
\end{array}%
\right] \left[ 
\begin{array}{cc}
0.944 & 0.056 \\ 
0.187 & 0.813%
\end{array}%
\right] =\left[ 
\begin{array}{cc}
0.901 & 0.099 \\ 
0.328 & 0.672%
\end{array}%
\right] 
\]%
\[
P^{8}=P^{4}P^{4}=\left[ 
\begin{array}{cc}
0.8450 & 0.1550 \\ 
0.5168 & 0.4832%
\end{array}%
\right] 
\]%
\[
P^{16}=P^{8}P^{8}=\left[ 
\begin{array}{cc}
0.7941 & 0.2059 \\ 
0.6864 & 0.3136%
\end{array}%
\right] 
\]%
\[
P^{32}=P^{16}P^{16}=\left[ 
\begin{array}{cc}
0.7719 & 0.2281 \\ 
0.7603 & 0.2397%
\end{array}%
\right] 
\]%
\[
P^{64}=P^{32}P^{32}=\left[ 
\begin{array}{cc}
0.7693 & 0.2307 \\ 
0.7691 & 0.2309%
\end{array}%
\right] 
\]%
\[
P^{128}=P^{64}P^{64}=\left[ 
\begin{array}{cc}
0.7692 & 0.2308 \\ 
0.7692 & 0.2308%
\end{array}%
\right] 
\]

%TCIMACRO{\TeXButton{EndFrame}{\end{frame}}}%
%BeginExpansion
\end{frame}%
%EndExpansion

%TCIMACRO{\TeXButton{BeginFrame}{\begin{frame}}}%
%BeginExpansion
\begin{frame}%
%EndExpansion

\QTR{frametitle}{Why the initial distribution does not matter}

\begin{itemize}
\item Notice the rows of $P^{128}$,%
\[
P^{128}=P^{64}P^{64}=\left[ 
\begin{array}{cc}
0.7692 & 0.2308 \\ 
0.7692 & 0.2308%
\end{array}%
\right] 
\]

\item They are identical

\item Let $p_{0}=\left[ 
\begin{array}{cc}
p_{01} & p_{02}%
\end{array}%
\right] $

\item Then, since $p_{01}+p_{02}=1$ 
\begin{eqnarray*}
p_{0}P^{128} &=&\left[ 
\begin{array}{cc}
p_{01} & p_{02}%
\end{array}%
\right] \left[ 
\begin{array}{cc}
0.7692 & 0.2308 \\ 
0.7692 & 0.2308%
\end{array}%
\right] \\
&=&\left[ 
\begin{array}{cc}
0.7692 & 0.2308%
\end{array}%
\right]
\end{eqnarray*}

\item Initial distribution does not matter in the long run

\begin{itemize}
\item for the unconditional distribution
\end{itemize}
\end{itemize}

%TCIMACRO{\TeXButton{EndFrame}{\end{frame}}}%
%BeginExpansion
\end{frame}%
%EndExpansion

%TCIMACRO{\TeXButton{BeginFrame}{\begin{frame}}}%
%BeginExpansion
\begin{frame}%
%EndExpansion

\QTR{frametitle}{Value functions with markov chains}

\begin{itemize}
\item The value functions for our economy can be written as%
\begin{eqnarray*}
V(k_{t},A_{1}) &=&\max_{k_{t+1}}\left[ u(A_{1}f(k_{t})+(1-\delta
)k_{t}-k_{t+1})\right. \\
&&\left. +\beta \left[ p_{11}V\left( k_{t+1},A_{1}\right) +p_{12}V\left(
k_{t+1},A_{2}\right) \right] \right] ,
\end{eqnarray*}%
and%
\begin{eqnarray*}
V(k_{t},A_{2}) &=&\max_{k_{t+1}}\left[ u(A_{2}f(k_{t})+(1-\delta
)k_{t}-k_{t+1})\right. \\
&&\left. +\beta \left[ p_{21}V\left( k_{t+1},A_{1}\right) +p_{22}V\left(
k_{t+1},A_{2}\right) \right] \right] ,
\end{eqnarray*}

\item Note the probabilities in each equation

\item These can be solved recursively

\begin{itemize}
\item beginning with some $V_{0}(\cdot ,A_{1})$ and $V_{0}(\cdot ,A_{2})$

\item just need to keep track of which probabilities to use
\end{itemize}
\end{itemize}

%TCIMACRO{\TeXButton{EndFrame}{\end{frame}}}%
%BeginExpansion
\end{frame}%
%EndExpansion

%TCIMACRO{\TeXButton{BeginFrame}{\begin{frame}}}%
%BeginExpansion
\begin{frame}%
%EndExpansion

\QTR{frametitle}{Example economy }

\begin{itemize}
\item We use a markov chain of 
\[
P=\left[ 
\begin{array}{cc}
.9 & .1 \\ 
.4 & .6%
\end{array}%
\right] 
\]

\item Get the plans of (similar but not identical to the last problem)\FRAME{%
ftbpFU}{3.039in}{2.2857in}{0pt}{\Qcb{The plans with Markov chains}}{}{policy
function sto2.jpg}{\special{language "Scientific Word";type
"GRAPHIC";maintain-aspect-ratio TRUE;display "USEDEF";valid_file "F";width
3.039in;height 2.2857in;depth 0pt;original-width 5.8332in;original-height
4.3751in;cropleft "0";croptop "1";cropright "1";cropbottom "0";filename
'C:/macro 2 2022/class 3/policy function sto2.jpg';file-properties "XNPEU";}}
\end{itemize}

%TCIMACRO{\TeXButton{EndFrame}{\end{frame}}}%
%BeginExpansion
\end{frame}%
%EndExpansion

%TCIMACRO{\TeXButton{BeginFrame}{\begin{frame}}}%
%BeginExpansion
\begin{frame}%
%EndExpansion

\QTR{frametitle}{Simulated economy with markov chain (same shocks as in
other)}

\FRAME{ftbpFU}{3.039in}{2.2857in}{0pt}{\Qcb{A simulation with Markov chains}%
}{}{storun2 exp.jpg}{\special{language "Scientific Word";type
"GRAPHIC";maintain-aspect-ratio TRUE;display "USEDEF";valid_file "F";width
3.039in;height 2.2857in;depth 0pt;original-width 5.8332in;original-height
4.3751in;cropleft "0";croptop "1";cropright "1";cropbottom "0";filename
'C:/macro 2 2022/class 3/storun2 exp.jpg';file-properties "XNPEU";}}

\begin{itemize}
\item Note the increased presistence
\end{itemize}

%TCIMACRO{\TeXButton{EndFrame}{\end{frame}}}%
%BeginExpansion
\end{frame}%
%EndExpansion

%TCIMACRO{\TeXButton{BeginFrame}{\begin{frame}}}%
%BeginExpansion
\begin{frame}%
%EndExpansion

\QTR{frametitle}{Computer program for Markov chains}

\texttt{global vlast1 vlast2 beta delta theta k0 kt At p1 p2}

\texttt{hold off}

\texttt{hold all}

\texttt{vlast1=20*ones(1,40);}

\texttt{vlast2=vlast1;}

\texttt{k0=0.4:0.4:16;}

\texttt{kt11=k0;}

\texttt{kt12=k0;}

\texttt{beta=.98;}

\texttt{delta=.1;}

\texttt{theta=.36;}

\texttt{A1=1.75;}

\texttt{p11=.9;}

\texttt{p12=1-p11;}

\texttt{p21=.4;}

\texttt{p22=1-p21;}

\texttt{A2=.75;}

%TCIMACRO{\TeXButton{EndFrame}{\end{frame}}}%
%BeginExpansion
\end{frame}%
%EndExpansion

%TCIMACRO{\TeXButton{BeginFrame}{\begin{frame}}}%
%BeginExpansion
\begin{frame}%
%EndExpansion

\texttt{numits=250;}

\texttt{for k=1:numits}

\texttt{\ \qquad for j=1:40}

\texttt{\ \qquad \qquad kt=k0(j);}

\texttt{\ \qquad \qquad At=A1;}

\qquad \qquad \texttt{p1=p11;}

\texttt{\ \qquad \qquad p2=p12;}

\texttt{\qquad \qquad\ z=fminbnd(@valfunsto,.41,15.99);}

\qquad \qquad \texttt{v1(j)=-valfunsto(z);}

\qquad \qquad \texttt{kt11(j)=z;}

\texttt{\ \qquad \qquad At=A2;}

\texttt{\ \qquad \qquad p1=p21;}

\texttt{\ \qquad \qquad p2=p22;}

\qquad \qquad \texttt{z=fminbnd(@valfunsto,.41,15.99);}

\texttt{\ \qquad \qquad v2(j)=-valfunsto(z);}

\texttt{\ \qquad \qquad kt12(j)=z;}

\texttt{\ \qquad end}

%TCIMACRO{\TeXButton{EndFrame}{\end{frame}}}%
%BeginExpansion
\end{frame}%
%EndExpansion

%TCIMACRO{\TeXButton{BeginFrame}{\begin{frame}}}%
%BeginExpansion
\begin{frame}%
%EndExpansion

\texttt{\ \qquad if k/50==round(k/50)}

\texttt{\ \qquad \qquad plot(k0,v1,k0,v2)}

\texttt{\ \qquad \qquad drawnow}

\texttt{\ \qquad end}

\texttt{\ \qquad vlast1=v1;}

\texttt{\ \qquad vlast2=v2;}

\texttt{end}

\texttt{hold off}

\texttt{\%plot(k0,kt11,k0,kt12)}

%TCIMACRO{\TeXButton{EndFrame}{\end{frame}}}%
%BeginExpansion
\end{frame}%
%EndExpansion

%TCIMACRO{\TeXButton{BeginFrame}{\begin{frame}}}%
%BeginExpansion
\begin{frame}%
%EndExpansion

\QTR{frametitle}{Subroutine valfunsto}

Note that interpolation of the previous value function is linear.

\texttt{function val=valfunsto2(x)}

\texttt{global vlast1 vlast2 beta delta theta k0 kt At p1 p2}

\texttt{k=x;}

\texttt{g1=interp1(k0,vlast1,k,'linear');}

\texttt{g2=interp1(k0,vlast2,k,'linear');}

\texttt{kk=At*kt\symbol{94}theta-k+(1-delta)*kt;}

\texttt{val=log(kk)+beta*(p1*g1+p2*g2);}

\texttt{val=-val;}

%TCIMACRO{\TeXButton{EndFrame}{\end{frame}}}%
%BeginExpansion
\end{frame}%
%EndExpansion

\end{document}
